\chapter{Самосогласованное очищение состояния мостом и страты частичных изометрий}

Все предыдущие проверки считали семейное состояние $R$ и мост $B$
независимыми переменными. Именно это позволило фиксировать чистое
$R=P$, одновременно отправляя поперечный блок $C$ в бесконечность.
Проверим каноническую альтернативу: мост является амплитудой очищения, а
состояние --- его нормированным редуцированным квадратом.

\section{Два редуцированных состояния одного моста}

Для $B\ne0$ определим
\begin{equation}
 R_R(B)=\frac{B^TB}{\Tr(B^TB)},
 \qquad
 R_L(B)=\frac{BB^T}{\Tr(B^TB)}.
 \label{eq:v6-self-consistent-reduced-states}
\end{equation}
Они имеют одинаковые ненулевые собственные значения и являются правым и
левым редуцированными состояниями вещественной амплитуды $B$. Такой
словарь стандартен для очищений и индуцированных мер на матрицах
плотности; см. \path{arXiv:quant-ph/0012101} и
\path{arXiv:1508.02279}.

Чтобы сохранить левую--правую симметрию исходного действия, положим
\begin{equation}
 \mathcal S_{\rm sc}(B)=\frac17\left[
 \Tr R_R(B)(I-B^TB)^2
 +\Tr R_L(B)(I-BB^T)^2
 \right].
 \label{eq:v6-self-consistent-bridge-action}
\end{equation}
Новый коэффициент не введён. При $B=\lambda I_3$ состояния равны
$I_3/3$, и точно восстанавливается прежнее скалярное действие
\begin{equation}
 \mathcal S_{\rm sc}(\lambda I_3)
 =\frac27(1-\lambda^2)^2.
\end{equation}

\section{Точная сингулярная формула}

Пусть $x_i=s_i^2$ --- квадраты сингулярных значений $B$, а
$T=x_1+x_2+x_3$. Тогда
\begin{equation}
 \boxed{
 \mathcal S_{\rm sc}(B)
 =\frac{2}{7T}\sum_{i=1}^3x_i(1-x_i)^2.}
 \label{eq:v6-self-consistent-singular-action}
\end{equation}
Она неотрицательна и при $\|B\|\to\infty$ растёт как $\|B\|^4$.
Следовательно, прежняя бесконечная долина исчезает, а матричный интеграл
становится коэрцитивным на бесконечности.

\begin{theorem}[Самосогласование устраняет расходимость]
При связи \eqref{eq:v6-self-consistent-reduced-states} условие
$R_R=P$ эквивалентно рангу $B=1$. Для семейства
$B=\operatorname{diag}(1,C)$ это требует $C=0$. Поэтому прежняя
четырёхмерная долина при фиксированном чистом состоянии не принадлежит
самосогласованному пространству конфигураций.
\end{theorem}

Это первый механизм, устраняющий убегание без логарифмического барьера,
обратного веса или пространственной производной.

\section{Нулевое множество действия}

Из \eqref{eq:v6-self-consistent-singular-action} следует
\begin{equation}
 \mathcal S_{\rm sc}(B)=0
 \quad\Longleftrightarrow\quad
 x_i\in\{0,1\},\qquad B\ne0.
\end{equation}
То есть нули --- все ненулевые частичные изометрии рангов
$k=1,2,3$.

Их редуцированные состояния имеют спектры
\begin{align}
 k=1:&\quad(1,0,0),\\
 k=2:&\quad(1/2,1/2,0),\\
 k=3:&\quad(1/3,1/3,1/3).
\end{align}
Ранг один действительно даёт чистую орбиту $\mathbb{RP}^2$ на правой
стороне. Однако то же действие с той же нулевой энергией допускает ранги
два и три.

\begin{proposition}[Коэрцитивность без выбора ранга]
Самосогласованное действие делает интеграл конечным, но не выбирает
материальную страту. Замкнутый ортогональный мост, рангово-два частичная
изометрия и чистый рангово-один мост являются вырожденными нулями.
\end{proposition}

\section{Энтропия выбирает замкнутый мост}

Добавим уже использовавшийся канонический член состояния
\begin{equation}
 \mathcal F_{\rm sc}(B)
 =\mathcal S_{\rm sc}(B)
 +\Tr\bigl(R_R(B)\log R_R(B)\bigr).
 \label{eq:v6-self-consistent-free-energy}
\end{equation}
Для любого qutrit-состояния
\begin{equation}
 \Tr(R\log R)\ge-\log3,
\end{equation}
причём равенство достигается только при $R=I_3/3$. Поскольку
$\mathcal S_{\rm sc}\ge0$, глобальный минимум
\eqref{eq:v6-self-consistent-free-energy} точно равен
\begin{equation}
 \mathcal F_{\min}=-\log3
\end{equation}
и достигается на $B\in O(3)$.

На трёх нулевых стратах значения равны
\begin{equation}
 0,\qquad -\log2,\qquad -\log3.
\end{equation}
Следовательно, энтропия строго выбирает полноранговый изотропный мост.

\begin{theorem}[Самосогласованное очищение не рождает проекторную фазу]
Связь $R=R_R(B)$ устраняет непертурбативную расходимость, но вместе с
канонической энтропией имеет единственный тип глобального минимума:
$B\in O(3)$ и $R=I_3/3$. Чистая орбита $\mathbb{RP}^2$ остаётся
допустимой нулевой стратой мостового действия, но не минимумом полной
свободной энергии.
\end{theorem}

\section{Коррекция статуса прежнего самозапуска}

Отрицательная однопетлевая мода $-45/16$ была получена интегрированием
флуктуаций $B$ при фиксированном независимом $R$. Если $R$ определяется
самим $B$, нельзя одновременно интегрировать $B$ и считать его
редуцированное состояние внешним фоном. Требуется обычное фоновое
разложение одного самосогласованного поля, а прежняя формула не переносится
автоматически.

Поэтому условный самозапуск не опровергнут арифметически, но его ключевая
предпосылка --- независимость $R$ и $B$ --- несовместима с механизмом,
который устраняет долину.

\begin{nogocriterion}[Запрет одновременной независимости и очищения]
Нельзя использовать $R(B)$ для удаления расходимости, сохраняя при этом
однопетлевой детерминант, выведенный для независимого фиксированного $R$.
Это два разных пространства конфигураций и два разных гессиана.
\end{nogocriterion}

\section{Вердикт}

\begin{protocol}
\begin{enumerate}
 \item канонические лево-правые редуцированные состояния:
 \textbf{построены};
 \item новый непрерывный коэффициент:
 \textbf{не введён};
 \item восстановление скалярного мостового действия:
 \textbf{точное};
 \item прежняя долина $B=\operatorname{diag}(1,C)$ при $R=P$:
 \textbf{устранена};
 \item рост действия на бесконечности:
 \textbf{квартичный};
 \item конечность матричного интеграла:
 \textbf{восстановлена};
 \item нулевые страты:
 \textbf{частичные изометрии рангов 1, 2 и 3};
 \item автоматический выбор ранга один:
 \textbf{нет};
 \item выбор энтропией:
 \textbf{ранг 3 и $R=I_3/3$};
 \item прежний независимый однопетлевой самозапуск:
 \textbf{не переносится};
 \item следующая проверка:
 \textbf{селектор ранговой страты частичной изометрии};
 \item рождение материи:
 \textbf{не закрыто}.
\end{enumerate}
\end{protocol}

Следующий тест должен проверить, существует ли в уже доказанной
топологии, Real-структуре или мере причина выбрать рангово-один страту
среди трёх нулевых страт. Если такой причины нет, ограничение
$\rank B=1$ будет эквивалентно предположению материи заранее.

Машинный сертификат сохранён в
\path{s2t_v6_self_consistent_state_bridge_purification_gate_results.json}.